\section{Let successful reads supply the message}
The file routine owns the open descriptor, a private hashing context and the caller's
selected scratch interval. Open the supplied terminated path read-only. If open fails,
return status one and leave the digest unchanged. Otherwise initialize the private
context and begin reading into scratch.

A positive read result is a byte count. Pass exactly that count to update, even if it
is smaller than the request, then read again. The scratch buffer may contain old bytes
after the returned prefix. They are not part of the message. A zero result is EOF:
finalize the context into the caller's output, close the descriptor and return status
zero. A negative result is a read failure: close the descriptor, return status two and
leave the output unchanged.

Preserving output on failure makes publication a deliberate boundary. The output may
still contain an earlier digest or a recognizable test pattern. The caller must read
the status before treating those bytes as this file's result. Do not finalize a
successfully read prefix after a later error and report it as the digest of the file.
The prefix and the complete file are different messages.

The descriptor must survive calls to the hash routines. Keep it in a preserved register
or private stack slot. The same is true of scratch, capacity and output pointers. A
nested \texttt{jal} also replaces \texttt{ra}; save the caller's return address before
making the first nested call. The private context occupies 368 bytes, separate from
saved registers, so its schedule writes cannot overwrite those saved values.

Use the same exact file for an external digest comparison. A text editor may add a
newline when saving, and a shell command that prints text may add one too. Begin with
fixtures whose length and bytes are known. Only after those agree should a mismatch
be treated as an arithmetic problem. The ordinary-file tests and controlled short-read
or read-error tests establish different parts of this interface.
